\documentclass[11pt]{article}
\usepackage{amsmath,amssymb,amsthm,geometry}\geometry{margin=1in}
\newtheorem{theorem}{Theorem}\newtheorem{proposition}[theorem]{Proposition}
\title{Shared-world feedback and external post-tightening}
\date{16 November 2026}
\begin{document}\maketitle
\paragraph{Audit and interface.} Shared-world compatibility theorem. It uses the deterministic port-free finite-feedback construction, whole-law root/block decomposition, and actual common-boundary limits. It proves compatibility of feedback with postcomposition by a continuation at external outputs, including simultaneous cross-invocation and continuation scopes. A common full primitive-diagram writer library is stated explicitly. It is not silently inferred from isomorphic realized scope bodies or from distinct free schema grammars. Scope writer and terminal polynomial grades remain separate.

\paragraph{Dependencies.} Shared-world assembly v3 has SHA256 \texttt{bbc0006ed160be930687306e2f57b6ab57d111e93b13e6c92b4f80dc881d60cd}; whole-scope decomposition v3 has \texttt{d3a9f7dfb58de84e685ec41983b06ed4914fa615d76b6aa23040627e328b1fb3}; reviewed free-global feedback v2 has \texttt{5568c1ae5781688c1dad959e6b8dd97d6685dfb43a660139d3b9a9dfc6e61390}. The source \emph{Native Post-Tightening Grade Comparison}, SHA256 \texttt{7ac6942aed9671053ee3190754d39b1acd19b2f46a9bdef2364acc5745a4e147}, was read directly for the uncolored full-grade theorem and its arity obstruction; its coefficient proof uses symmetric monoidal coherence, independently of the operational world law. The new claim here is shared-world whole-scope compatibility, not novelty of that terminal-grade map.

\section{Inputs and the common primitive-diagram library}
Let $P:A+U\rightsquigarrow B+U$ have deterministic port-free seeds. Recursive invocations of its $U$ branch satisfy the reviewed finite-carrier interface: persistent ordered ports, actual world-arrow rows, passive rebase, completed-body stability and finite active births. Let $Q:B\rightsquigarrow Z$ satisfy the same local operational requirements except that its outputs are all external $Z$ outputs; it has no recursive $U$ outputs. Its initial configuration is finite and can already emit finitely many $Z$ ports. All states, including unreachable ones, are retained. Tags distinguish $A,U,B,Z$ branches even if underlying sets intersect.

Use the common-boundary limit mode with all admitted whole joint laws, the same fixed world category and full boundary library, and all-raw compatible primitive body cones. Let $L_Q$ be the typed continuation which uses $Q$ at $B$ and a pure zero-actor identity at $U$. An identity contributes no command, paid step, sample or context/writer factor. Define the operational composite $R=P;L_Q$ by immediate continuation initialization at emitted ports, not waiting for completion. It has port-free seeds because $P$ does.

Independently declare a full library $\mathscr D_{P,Q}$ of finite typed addressed active primitive diagrams over the fixed boundary library. A primitive position records either a $P$ invocation address or a $Q$ copy attached to a particular persistent $B$ port of a $P$ address. Stable port keys qualified by their frozen input labels, not changing ranks, determine these addresses; pure-case tags are removed only by the specified fixed identity-wire normalization; active planar order is retained separately. A diagram includes the primitive schema labels, boundary incidences and limit mode. Include all declared formal diagrams and their full limits, whether or not some configuration realizes them. In particular do not define the library as the union of positive-support executions. An implementation may take every finite well-typed such diagram, including ones that cannot jointly occur in a configuration. Each family is the actual full-family limit of its primitive factors. Every unused shape and coefficient is retained.

Both operational systems below are equipped with this same library and label each actual scope by its flat active diagram. This is a direct instance of the source's full primitive/joint-library interface. If instead starting from separately generated free schema libraries, use their global flattening maps into $\mathscr D_{P,Q}$ and the actual limit comparisons on every formal schema. These are native interpretation maps; they are not asserted to be invertible or to preserve initiality after adding or identifying unused schema names. A full common-library declaration is a substantive index choice, not an unnoticed quotient by equal observed bodies.

\paragraph{Global flattening on unrealized schemas.} The left generated schema grammar has an optional active feedback-$P$ schema and active/passive $Q$ names at the external-port positions of its root frame. Flatten the former at all its primitive $P$ addresses and prefix the latter by their tagged owning $B$-port addresses. The right grammar has an invocation plan whose active primitive names are themselves the formal $P;L_Q$ joint schemas. Expand each such name into its declared primitive $P/Q$ factors, prefix by its invocation address, and remove only the fixed pure-identity tags. These formulas apply to all formal names, not just enabled commands. Labels and boundary arrows must be well typed before flattening; an ill-typed diagram is not repaired by forgetting its type. A diagram may be dynamically unrealizable while still being a well-typed finite diagram of full primitive families. The common library admits those diagrams independently. At each source name the full-family comparison is the actual iterated-to-flat limit isomorphism, determined by every typed primitive projection. Thus all unused shapes and actual coefficient arrows are mapped. When two distinct generated names flatten to one diagram, that name identification is the declared interpretation kernel; it is not a claim that the original libraries were isomorphic. All shape-indexed libraries are small in an adequate universe and the ambient family category admits all such small families; only finite coefficient limits are used.

\section{Actual state and whole-scope comparison}
Write $\mathrm{Fb}_U$ for the finite invocation feedback construction with deterministic port-free seeds. The left system
\[
 X=\mathrm{Fb}_U(P);Q
\]
has one finite $P$-invocation tree, and a $Q$ configuration at every currently emitted external $B$ port anywhere in that tree. These $Q$ copies begin at the birth of their ports and can act while ancestors or other $P$ invocations remain unfinished. The right system
\[
 Y=\mathrm{Fb}_U(P;L_Q)
\]
has a tree of $R$ invocations. Each local $R$ state contains one $P$ state and the $Q$ copies at its own $B$ ports; the emitted $U$ ports immediately start child $R$ invocations. Remove the fixed pure identity tags on $U$ wires and retain all other stable path keys. This is an invertible global reparse of exactly the same primitive states, shared world, ports and incidence. It does not merge two occurrences.

\begin{theorem}[Whole-scope operational post-tightening]
The reparse is an isomorphism of the complete shared-world operational interfaces of $X$ and $Y$: configurations, initialization, rebase, external ports, completion, menus, all raw scope atoms and laws, actual context coefficients, full scope writer-library maps and bodies. The terminal-grade comparison is treated separately below.
\end{theorem}
\begin{proof}
A command on either side selects a nonempty set of existing $P$ and $Q$ primitive actors, one local command at each, and one admitted joint law on their full matching raw diagram. On the left it is presented by the $P$-feedback block and the external $Q$ blocks; on the right by the local $R$ invocation blocks, each containing its $P$ actor and local $Q$ copies. Starting with the full primitive joint, restrict it to each relevant block. Nonempty blocks carry the uniquely induced marginal joint; empty blocks carry the response to the same actual boundary witness. Retain the outer joint as the full pushed flat law. Conversely these block data reconstruct the flat law by the tuple bijection. All-joint admission makes each marginal legal. This gives inverse menu/row translations. It neither takes a product of marginal laws nor serializes a joint event, and preserves zero tuples and boundary values of zero mass.

The actual context coefficient comparison is the canonical comparison between two iterated limits of the same primitive coefficient diagram at each raw atom. The full writer map is identity on the already common flat library. If a grouped presentation still uses iterated products, its comparison to this library is the canonical actual full-family limit comparison on every formal schema, not only on its realized body. Coherence follows by testing each primitive and boundary projection. Pure identity factors are absent; passive actor blocks contribute only the retained slice boundary.

For posts, every old selected $P$ or $Q$ actor has the same active post under the same raw primitive coordinate and actual world arrow. Every old unselected actor has the same passive rebase. A new $B$ port creates a $Q$ seed at the actual post-world on both sides. A new $U$ port creates a $P$ seed on the left and an $R$ seed on the right. These agree after reparse because the $P$ seed is port-free and therefore the $R$ seed initially has no $Q$ continuation. No old child is reinitialized. No actor born during the scope is selected in that scope. Existing and fresh external $Z$ ports have the same stable order and labels. These checks apply to every raw atom and unreachable configuration. Initial states are the same port-free $P$ seeds; completion on both sides means all primitive $P$ and $Q$ actors are complete. Rebase commutes by its componentwise definition. Hence all listed native fields commute.
\end{proof}

This theorem includes a scope simultaneously involving a $P$ ancestor, a $P$ descendant and a $Q$ continuation attached to a different invocation. Their joint law may be correlated and the primitive commands need not commute. Block transport uses one global law and one common update; separate executions with an independence assumption are not substituted.

\section{Full terminal grades: colored isomorphism versus uncolored direction}
First use a genuinely colored terminal collection for $P$: each formal slot, including unused ones, is independently typed as an external $B$ slot or a recursive $U$ slot. Substitute the entire terminal collection $D$ of $Q$ only into external slots. Then a formal terminal shape on the left is a finite colored $P$ invocation tree with a $D$ shape at every external slot. The right shape stores the very same $D$ decorations locally at their owning $P$ nodes. Thus there is a bijection on \emph{all} formal shapes. It has the same ordered external positions and the same labeled primitive coefficient factors, regrouped by symmetric monoidal associators, unitors and permutations. These are actual inverse coefficient arrows. Together with the operational theorem this gives a full native isomorphism in this colored-grade doctrine.

The source's original whole constant-polynomial case construction is different. With whole grades $C,D$ it retains the full case grade $I+D$ even when dynamic typing selects only one branch. Its entire feedback grades are
\[
 (R_UC)\circ D\quad\hbox{and}\quad R_U(C\circ(I+D)).
\]
The source's reviewed map sends each original internal invocation edge to an $I$ decoration and each original external slot to its $D$ decoration. It is injective on full shapes, bijective on ordered external positions, and invertible on each actual coefficient, but it is not generally surjective on shapes. Every unused source shape is mapped. Combined with the shared-world operational theorem it gives a directed full native post-tightening comparison, not a native isomorphism.

For the exact reverse obstruction take one loop label, unary $C$ and binary $D$, with unit coefficients. Every shape of $(R_UC)\circ D$ has arity two. The right collection contains a root decorated by $D$ with one output external and one output continuing to another $D$-decorated node; it has arity three. No arity-preserving full grade arrow can map that right collection into the left. The impossible shape is not realized by typed operational execution but remains part of the declared full uncolored grade. Neither shared-world operational isomorphism nor writer-library flattening repairs this obstruction.

\section{Nontrivial full-coefficient test}
Take augmented rational vector spaces for both context and the independent writer coefficients, with boundary coefficient the augmentation unit. At three active actors, for example a $P$ ancestor, its $P$ child and a $Q$ continuation attached to the ancestor, let each raw row be a fair bit with the same actual update arrow $w\to w+1$. Retain the full eight-tuple carrier and the diagonal joint of mass $1/2$ on $(0,0,0)$ and $(1,1,1)$ and zero on the other six tuples. Active states cache their respective bits; passive actors retain their own state by rebase. Persistent ports witnessing these three actors can be supplied in an admitted configuration, independently of whether that configuration is reachable from the seed.

Give each primitive context atom $K=(\mathbb Q^2,\operatorname{pr}_1)$. Its flat three-actor coefficient is $K\times_{\mathbb Q}K\times_{\mathbb Q}K$, of dimension four. Both block decompositions produce that same actual coefficient by canonical projection-preserving isomorphisms. The eight-dimensional tensor $K^{\otimes3}$ is not substituted for it.

Independently give each active primitive writer schema two formal shapes, with coefficient $K$ on the first and $M=(\mathbb Q^3,\operatorname{pr}_1)$ on the second, both over the singleton unit boundary family. Let actual bodies use the first shape. The full three-actor writer family nevertheless has all eight shape triples. At a triple with $r$ second-shape choices the actual limit coefficient has dimension $4+r$. The block comparison is an actual isomorphism on each of these coefficients, including all seven unrealized shape triples. It is forced by the same labeled factor projections, not by the single realized body. This tests the symbolic full-family limit argument with nontrivial unused coefficients while retaining the independent terminal grade doctrine.

\section{Coherence and comparison scope}
For further continuation $S:Z\rightsquigarrow V$, the operational routes obtained by first attaching $Q$ then $S$, or attaching the composite continuation, flatten to the same addressed family of primitive $P,Q,S$ actors. The global scope law is the same full law on that family. Actual context and full writer maps agree by primitive projections; post-state maps agree actorwise and initialize each newborn continuation at the same actual post-world. This proves coherence in a common full $P,Q,S$ diagram library without a serializability assumption.

For colored terminal collections the factor/regrouping proof gives the corresponding isomorphism coherence. For the original uncolored collections the common-target comparison uses the source map
\[
 I+(D\circ E)\longrightarrow (I+D)\circ(I+E),
\]
which chooses $I$ then $I$ on the internal branch and $D$ then $E$ at every external position on the external branch. This is directed because unused mixed alternatives remain on the target. The two tightening routes agree on every original full shape and every labeled coefficient factor; they are not asserted to identify all intermediate full grade collections.

Naturality is claimed for native parameter comparisons that supply coherent common-library transport on all primitive schemas and stable occurrence addresses, in addition to their local port, state and row maps. Under those hypotheses both routes apply the same primitive maps, push the same whole joint and use the same induced actual limits. A merely state-dependent port renaming with no global library action is not silently promoted to a full-library natural transformation.

\paragraph{Contract consequence and residual.} This is a direct feedback/open-substitution compatibility law in the genuinely shared-world interface, not two independent operations placed side by side. Its full native strength depends separately on the writer-library choice and on the colored versus uncolored terminal doctrine. The free/flat feedback initiality statements retain their previously specified algebra categories; equipping a runtime with a larger common unused library does not automatically preserve their universal property. Strong original-source inclusion, random finite initialization, arbitrary cyclic feedback and a single merged content grade remain outside this result.
\end{document}
